\section{Results}
\vspace{1cm}
We compare five generative priors against two LASSO baselines. Four of the priors are the models described in the previous section, a VAE and a DCGAN each trained with a latent vector of dimension 20 and 30. The fifth is the fully connected architecture of Bora et al. \cite{bora2017}, included so that the comparison with the numbers they publish is direct rather than indirect.

\subsection{Experimental Protocol}

Every method is evaluated on the same ten test images, one for each digit class, so that the comparison is not driven by a single lucky or unlucky example. At each measurement budget \( m \) all the methods receive the same measurement matrix \( A \) and the same noise vector, and the noise is scaled so that its expected norm remains equal to 0.1 whatever the value of \( m \).
\\
Ten is a small evaluation set, and the reason is where the cost of this experiment sits. Recovery with a generative prior is not a forward pass through the network but an optimization in its latent space, a thousand gradient steps from ten random initializations, carried out for every image at every budget. A DCGAN generator costs about ninety times what the cheapest prior costs and several hundred times what LASSO costs, whatever machine it runs on. That cost is linear in the number of images and is paid again for every method, every budget and every setting of the penalty, so a fixed computing budget buys either a larger evaluation set or a wider comparison. We chose the comparison: seven methods at ten budgets, the whole sweep repeated without the regularization term, and a further sweep over four values of it. Stratifying by class matters more than the count at this size, since ten images drawn at random can miss a digit class or repeat one, and the methods do not all fail on the same digits. What the size of this set costs in statistical power is reported in section 6.5, and enlarging it is the first of the extensions discussed in the conclusions.
\\
The reconstruction quality is measured as the squared distance between the reconstruction and the original image, divided by the number of pixels:
\begin{equation}
\text{error per pixel} = \frac{\lVert \hat{x} - x^* \rVert_2^2}{n}, \qquad n = 784 .
\end{equation}
This is not the same quantity as the measurement error \( \lVert AG(\hat{z}) - y \rVert \) that the recovery algorithm minimizes. The latter decreases as \( m \) decreases, simply because fewer constraints have to be satisfied, so it describes how well the optimization converged rather than how good the reconstruction is. For this reason it is used only as a diagnostic and to rank the random restarts, while every number reported below is computed against the ground truth image.
\\
All the reconstructions obtained with a generative prior use a thousand gradient steps, ten random restarts and the regularization weight \( \lambda = 0.1 \) recommended in \cite{bora2017}. Section 6.6 reports what happens without the regularization term.

\subsection{The LASSO Baselines}

LASSO (Least Absolute Shrinkage and Selection Operator) assumes sparsity in some basis \( \Psi \), writes \( x = \Psi\theta \), and solves
\begin{equation}
\hat{\theta} = \arg \min_\theta \|y - A\Psi\theta\|_2^2 + \alpha \|\theta\|_1, \qquad \hat{x} = \Psi\hat{\theta} .
\end{equation}
\\
We report two choices of \( \Psi \). The first is the pixel basis, \( \Psi = I \), which is the baseline used in \cite{bora2017} for MNIST: the authors observe that the digits are already sparse in pixel space, since most of the image is background. The second is the orthonormal DCT basis, which the same authors use for natural images and which is the usual choice for image compression. The two behave very differently on this dataset and reporting only one of them would misrepresent how strong the classical approach is.
\\
A baseline is only informative if it is given a good configuration, so the shrinkage \( \alpha \) was swept over eight values spanning seven orders of magnitude, separately for each basis and for each measurement budget, and the value minimizing the error was kept at every budget. Tuning per budget rather than once is deliberate: the implementation we use scales the objective by \( 1/(2m) \), so a fixed \( \alpha \) is effectively a much stronger penalty at large \( m \) than at small \( m \), and a single value costs the baseline up to a factor of two at the ends of the range. The sweep is run on ten images drawn with a different seed from the ones the benchmark scores, and with its own measurement matrices and noise, so the baseline is given its best configuration without that configuration being chosen on the evaluation data. The reconstructions are clipped to \( [0,1] \), the range MNIST pixels live in, which also improves the baseline. Note that the shrinkage quoted in \cite{bora2017} is not directly transferable, for the same scaling reason.
\\
From our experiments, LASSO in either basis is inadequate until a substantial number of measurements is available, and neither produces a coherent digit at 200 measurements or below. The two fail differently: the DCT basis returns a noisy image in which the digit is buried, while the L1 penalty in the pixel basis drives most of the reconstruction to zero, so it returns an almost blank image with a few scattered dots. On the other hand LASSO keeps improving as measurements are added, and in the pixel basis it becomes essentially exact once the budget approaches the dimension of the signal.

\subsection{Comparison of the Methods}

Table \ref{tab:results} and figure \ref{fig:error vs measurements} report the error per pixel of every method as the number of measurements grows.

\begin{table}[H]
    \centering
    \small
    \begin{tabular}{r|cc|ccccc}
    \( m \) & LASSO & LASSO & VAE & VAE & VAE & DCGAN & DCGAN \\
     & (pixel) & (DCT) & (paper) & \( k{=}20 \) & \( k{=}30 \) & \( k{=}20 \) & \( k{=}30 \) \\
    \hline
    10  & 0.1217 & 0.1032 & \textbf{0.0680} & 0.0778 & 0.0755 & 0.0864 & 0.1060 \\
    25  & 0.1189 & 0.0986 & 0.0276 & 0.0288 & \textbf{0.0194} & 0.0813 & 0.0635 \\
    50  & 0.1171 & 0.0835 & \textbf{0.0118} & 0.0155 & 0.0191 & 0.0327 & 0.0383 \\
    75  & 0.1127 & 0.0710 & \textbf{0.0084} & 0.0088 & 0.0091 & 0.0167 & 0.0306 \\
    100 & 0.1046 & 0.0628 & 0.0083 & 0.0083 & \textbf{0.0077} & 0.0142 & 0.0147 \\
    200 & 0.0827 & 0.0325 & 0.0075 & 0.0079 & \textbf{0.0072} & 0.0125 & 0.0155 \\
    300 & 0.0408 & 0.0202 & 0.0072 & 0.0076 & \textbf{0.0071} & 0.0125 & 0.0102 \\
    400 & 0.0107 & 0.0115 & 0.0072 & 0.0076 & \textbf{0.0070} & 0.0120 & 0.0088 \\
    500 & \textbf{0.0008} & 0.0064 & 0.0072 & 0.0076 & 0.0071 & 0.0117 & 0.0096 \\
    750 & \textbf{0.0000} & 0.0006 & 0.0070 & 0.0075 & 0.0069 & 0.0116 & 0.0085 \\
    \end{tabular}
    \caption{Error per pixel, averaged over ten test digits, one per class. The best method at each budget is in bold. VAE (paper) is the fully connected architecture of \cite{bora2017}. Every generator was trained on the same 54,000 images, with the test split these digits come from held out.}
    \label{tab:results}
\end{table}

\begin{figure}[H]
    \centering
    \includegraphics[width=1\linewidth]{assets/error_vs_measurements.png}
    \caption{Error per pixel as a function of the number of measurements \( m \). Bars are 95\% percentile bootstrap intervals for the mean over the ten test images, from ten thousand resamples. \cite{bora2017} plot normal-theory intervals in their figure 1; a symmetric interval assumes the mean of ten strongly skewed errors is approximately normal, and at three points it puts the lower end below zero, outside the range a squared error can take. The bootstrap interval stays inside that range and shows the asymmetry. The significance tests in section 6.5 do not assume normality either; the bars are an illustration of spread and the tests are the inference. Both axes are logarithmic.}
    \label{fig:error vs measurements}
\end{figure}

Three regimes can be distinguished.
\\\\
\textbf{Few measurements.} Every VAE prior is more accurate than either baseline up to four hundred measurements. The DCGAN with latent dimension 30 is more accurate than both from twenty-five measurements upwards, the DCT baseline edging it at ten, and the one with latent dimension 20 is more accurate than both from ten to three hundred. These comparisons are against the baselines: among the priors themselves the two DCGANs remain the least accurate at every budget, as section 6.5 reports. The best prior leads by a wide margin throughout this range. With 25 measurements, the budget \cite{bora2017} single out for their own comparison, the best model has an error of 0.0194 against 0.0986 for LASSO in the DCT basis, a factor 5.1. The factor keeps growing with the budget in this range, so this is a conservative point at which to quote it. The ratio against LASSO in the pixel basis is 6.1, but that baseline deserves a caveat at this budget: predicting an all-zero image gives an error of 0.1178 on these digits, and the pixel baseline is at or above that level with 10 and 25 measurements, so it is not reconstructing anything there and a ratio against it carries little meaning. At that budget neither LASSO reconstruction is recognizable as a digit. This is the regime that motivates compressed sensing, and it is where the choice of the prior matters most.
\\\\
\textbf{Sample efficiency.} We measure it against each of the two baselines at 400 measurements, where they reach errors of 0.0115 in the DCT basis and 0.0107 in the pixel basis. All three VAE priors reach that level with 75 measurements, a saving of a factor \textbf{5.3}, and the factor is the same against either baseline, so the headline figure does not depend on which one we take as the reference. What does depend on that choice is the number of individual images beaten at that budget: 8 or 9 out of 10 against the DCT baseline and 1 out of 10 against the pixel one, for the reason given in the next paragraph. The DCGAN with latent dimension 30 needs 300 measurements, a factor 1.3; with latent dimension 20 it never reaches the level within our sweep. Bora et al. report a saving between 5 and 10 times: our reproduction falls at the lower end of that interval, and only the VAE priors fall inside it at all.
\\\\
The two baselines differ in kind at this budget, and the reason is worth stating. At 400 measurements the pixel baseline has a mean error of 0.0107 but a median of 0.0003: six of the ten digits are already recovered to better than \( 10^{-3} \) and one is still at 0.083. That budget falls in the middle of its transition from failure to near-exact recovery, so its mean describes the digits it has not solved yet rather than a level a method can be compared against, and the count of images beaten there is correspondingly low. The DCT baseline at the same budget has a mean of 0.0115, a median of 0.0113 and a worst case of 0.0151, so a per-image comparison against it means what it appears to mean. We report both rather than choosing one, since selecting the more favourable reference after seeing the answer would not be a measurement.
\\\\
\textbf{Many measurements.} From 500 measurements onwards LASSO in the pixel basis becomes the most accurate method by a large margin, reaching an error below \( 10^{-4} \) with 750 measurements against 0.0069 for the best generative model. Once the number of measurements approaches the 784 dimensions of the signal, an assumption of sparsity in pixel space is enough to recover an MNIST digit almost exactly, while the generative prior remains stuck. The reversal itself is expected and is described in \cite{bora2017}, where the authors observe that LASSO eventually surpasses the generative approach and that this requires more than 500 measurements. We observe it from 500 onwards, and in the DCT basis, which is the weaker baseline in this regime, from the same budget.
\\\\
The reason for the reversal is that the reconstruction is constrained to lie in the range of the generator. The distance between a real digit and that range does not depend on the number of measurements, so once enough measurements are available to identify the closest point in the range, additional measurements bring no benefit. Averaging the budgets from 300 upwards, this residual error settles at 0.0070 for the VAE with latent dimension 30, 0.0072 for the paper architecture, 0.0076 for the VAE with latent dimension 20, 0.0093 for the DCGAN with latent dimension 30 and 0.0119 for the DCGAN with latent dimension 20. These five values are properties of the generators and not of the recovery algorithm, and they span a factor of only 1.7, so on this side of the plot the generator matters far less than the gap to the baselines does at low budgets.

\subsection{Reconstruction Examples}

Figure \ref{fig:reconstruction grid} shows the same digit reconstructed by every method at every budget.

\begin{figure}[H]
    \centering
    \includegraphics[width=1\linewidth]{assets/reconstruction_grid.png}
    \caption{One test digit reconstructed by each method as the number of measurements increases}
    \label{fig:reconstruction grid}
\end{figure}

The visual comparison confirms the numbers. The DCT baseline returns a noisy image until about 300 measurements and the pixel baseline an almost blank one until about 400, while the generative priors produce a plausible digit almost immediately. The DCGAN gives visibly sharper strokes than the VAE, which is consistent with the fact that an adversarial generator is trained to produce samples a discriminator accepts, rather than an average of the plausible reconstructions as the VAE does. Sharpness and accuracy are not the same thing, however, and as the table shows the sharper model is the less accurate one.

\subsection{Architecture and Latent Dimension}

Because every method sees the same images and the same measurement matrices, the comparisons below are paired. We test them with a Wilcoxon signed-rank test on the ten per-image differences at each budget, preferring it to a paired t-test because the per-image errors are strongly skewed, a few digits being reconstructed far worse than the rest, and with ten samples the normality assumption of the t-test would be doing real work. Each pair of methods is tested at all ten budgets, so we correct the resulting p-values within that family with the Holm step-down procedure: it sorts the p-values, multiplies the smallest by ten, the next by nine and so on, and stops at the first that fails, which bounds the probability of even one false positive across the ten without assuming the tests are independent. Every p-value quoted below is the corrected one.
\\\\
\textbf{The generative priors beat sparse recovery, and then lose to it.} At 750 measurements every one of the ten prior and baseline pairs puts the baseline ahead, significantly in each case. In the other direction the joint statement is narrower than the means suggest: 50, 75, 100 and 300 are the only budgets at which every prior beats every baseline significantly at once, because the two DCGAN columns drop out elsewhere. Taken one at a time, each VAE prior is significantly better than both baselines from 25 to 300 measurements. The largest corrected p-value among these is 0.0488, and it belongs to a DCGAN column; every claim involving only the VAE priors sits at 0.0195, which is the smallest value attainable, since with ten paired samples the Wilcoxon test cannot go below 0.00195 and the correction across ten budgets multiplies that by ten.
\\\\
The correction is applied within each pair of methods, across the ten budgets, which is the family a claim such as "significant from 75 measurements up" spans. That choice is load bearing and not a formality: the table holds 200 tests in 20 such families, and a correction spanning all 200 would put the attainable minimum at 0.39, under which nothing in this report would be significant. We report the narrowest defensible family and say so rather than leaving the reader to assume a global correction.
\\\\
These tests are conditional on a single measurement matrix per budget. The pairing across images is genuine, but ten images observed through one draw of \( A \) are not ten independent sensing realizations, so the confidence intervals and the p-values describe variation between images and say nothing about variation between measurement matrices.
\\\\
\textbf{The three VAE priors are indistinguishable from one another.} No comparison among the fully connected architecture of \cite{bora2017}, our convolutional VAE with \( k = 20 \) and our convolutional VAE with \( k = 30 \) reaches significance at any budget: every corrected p-value is 1. Their error floors differ by 0.0005, which ten images are not enough to separate. We had expected to be able to rank the architectures, and to say something about the effect of the latent dimension, and the data does not support either statement. Any ordering read off the table would be noise, and we report the absence of an effect rather than the ordering.
\\\\
\textbf{The VAE beats the DCGAN at \( k = 20 \) everywhere, and at \( k = 30 \) only while measurements are scarce.} Against the DCGAN with latent dimension 20 every VAE is significantly more accurate at nine of the ten budgets. Against the one with latent dimension 30 the advantage is significant only while measurements are scarce, and at how many budgets depends on the VAE: four for our \( k = 30 \), three for the architecture of the paper and one for our \( k = 20 \). Elsewhere the difference is real but too small for ten images to establish: with 750 measurements that generator sits at 0.0085 against 0.0069 for the best VAE. Every VAE has a lower mean error than every DCGAN at every budget in the sweep; what changes with latent dimension 30 is only whether ten images suffice to prove it.
\\\\
\textbf{The latent penalty is not what separates the two families.} The table uses \( \lambda = 0.1 \), the value \cite{bora2017} report for their MNIST VAE, whereas the only value they report for a DCGAN is 0.001, obtained on celebA at a latent dimension of 100 and a different pixel range. Transferring it to a DCGAN on MNIST is our reading rather than their prescription, and \( \|z\|^2 \) scales with the latent dimension, so we ran both DCGANs at 0.001 as well rather than argue about it. The floor at \( k = 20 \) is essentially unchanged, 0.0116 against 0.0119, and the one at \( k = 30 \) is worse, 0.0104 against 0.0093. The penalty the main table uses is therefore not handicapping the DCGANs, and for one of them it is the better setting.

\subsection{The Latent Regularizer}

Figure \ref{fig:regularisation} compares the results obtained with \( \lambda = 0.1 \) and with the term removed.

\begin{figure}[H]
    \centering
    \includegraphics[width=1\linewidth]{assets/regularisation_comparison.png}
    \caption{Recovery with and without the latent regularization term}
    \label{fig:regularisation}
\end{figure}

The term behaves as its motivation suggests. It helps when the measurements are few, where the data alone does not determine the latent vector and the prior supplies the missing information: with 10 measurements the error of the best model falls from 0.0782 to 0.0680, an improvement of 13\%. It is mildly harmful once measurements are plentiful, where the same pull towards the prior keeps the solution from reaching the closest point in the range of the generator: with 750 measurements the error rises from 0.0054 to 0.0069.
\\
A sweep over \( \lambda \in \{0, 0.01, 0.1, 1\} \) makes the mechanism explicit. The norm of the recovered latent vector falls monotonically with \( \lambda \), from 9.95 to 2.67 on average, so the term does constrain the search as intended. More interestingly, the value of \( \lambda \) that minimizes the error is not fixed. It is 1 for all three priors at 10 measurements, and 0 for all of them at every budget of 200 and above; in between the three disagree. The single value recommended in \cite{bora2017} is therefore a compromise across regimes rather than an optimum at any one of them, which is consistent with the interpretation of the term as a substitute for missing measurements. It also means the generative side of our comparison runs a knowingly suboptimal regularizer at every budget, which works against it rather than for it. The sweep is computed on its own draw of images, matrices and noise, disjoint from the ones the benchmark scores, so that reading a preferred value off this table is not selection on the evaluation set. It covers the three VAE priors only, since recovery with a DCGAN costs tens of times more than with a VAE and the sweep repeats every recovery at four values of \( \lambda \).

\subsection{Computational Cost}

The methods are not comparable in cost. Reconstructing the ten test images at one budget, with ten restarts and a thousand gradient steps, takes on average about 1.5 seconds with either LASSO baseline, 6.6 seconds with the fully connected decoder of the paper, 15 and 19 seconds with our convolutional decoders and about 600 seconds with a DCGAN generator, a factor 93 between the cheapest and the most expensive generative prior. These are wall-clock times on a single CPU machine and are quoted for scale; the comparison between the methods is the ratio, which does not depend on the hardware.
\\
The parameter counts do not explain that ratio: the DCGAN generator has 2.9 million parameters against 0.65 million for the fully connected decoder, a factor of only 4.5, and our convolutional decoder has fewer parameters than the fully connected one yet is more than twice as slow. What the cost tracks is the amount of arithmetic per forward pass. The fully connected decoder is three matrix products, whereas the DCGAN generator applies transposed convolutions with 256 and then 512 channels at close to full resolution, so each of its ten thousand forward and backward passes moves orders of magnitude more data.
\\
This cost is incurred at every reconstruction and not only once during training, because recovery is itself an optimization problem. It is worth stressing that the two most expensive priors are also the two least accurate ones, at every budget, so on this dataset there is no trade-off to arbitrate between the two. The sweep reported in this section required about 3.5 hours of computation and the unregularized sweep of section 6.6 a comparable amount, the DCGAN recoveries accounting for almost all of both.

\subsection{Reproducibility of the Training}

While preparing these results we found that training the same VAE with the same script and different random seeds produced generators of very different quality. Measured as the best reconstruction reachable inside the range of the generator, the results ranged from 0.0103 to 0.0358 for \( k = 30 \), a factor of 3.5, which is comparable to the entire spread between the five priors we set out to compare.
\\
The cause is partial posterior collapse: in the weaker runs the decoder relies on a few latent directions and ignores the rest. Measuring how much each latent coordinate moves the generated image, the ratio between the average sensitivity and the largest one is 0.31 for the best model and 0.05 for the worst, and that ratio orders the models exactly as the reconstruction error does.
\\
Adding a KL warm-up, as described in section 3.2.3, addresses the problem. Every model improves, the best representation error at \( k = 30 \) falling from 0.0103 to 0.0062, and the KL divergence at convergence rises from between 4 and 7 nats to between 16 and 23. Across the two seeds of each configuration the remaining spread in representation error is a factor of 2.5, 1.3 and 1.2. That figure and the factor 3.5 above are not the same measurement, being taken on different images and over a different number of seeds, so we give them side by side rather than as a single ratio that shrank. The VAE checkpoints used in this report were selected on validation loss alone, and the DCGAN ones on representation error, a generative adversarial network having no validation loss to select on; both procedures were fixed before the runs were carried out.
\\
We report this because it affects how the comparisons above should be read: a difference between two priors is only meaningful when it is larger than the variability of the training procedure that produced them, which is why every comparison in section 6.5 is supported by a paired test rather than by a single pair of runs, and why the comparisons among the three VAE priors are reported as inconclusive rather than resolved.
